\documentclass[11pt]{article}
\usepackage{fontspec}
\usepackage[french,english]{babel}
\usepackage{amsmath}
\usepackage{amssymb}
\usepackage{amsthm}
\usepackage{mathrsfs}
\usepackage{tikz-cd}
\usepackage{enumitem}
\usetikzlibrary{babel}
\frenchsetup{AutoSpacePunctuation=false}

\newcommand{\printedpage}[1]{\marginpar{\footnotesize p.~#1}}
\newcommand{\msp}[1]{\marginpar{\footnotesize [#1]}}
\newcommand{\Hot}{\mathrm{Hot}}
\newcommand{\Hotz}{\mathrm{Hot}_0^{\wedge}}
\newcommand{\Hc}{\underline{\mathrm{Hom\,cof}}}
\newcommand{\Nop}{\mathbf{N}^o}

\begin{document}

\printedpage{74}
% head: A. Grothendieck : Dérivateurs, Ch. VII, Édition des Universités de Montpellier II et Paris VII
conditions at the beginning of the paragraph, $C_\infty$ must be written $C_\infty(\underline{t})$, it depends on the choice of the cofinal functor
\[
\underline{t} : \Nop \longrightarrow T.
\]
When $\underline{t}$ varies in the category
\[
\Hc(\Nop, T),
\]
the $C_\infty(\underline{t})$ vary contravariantly in $\underline{t}$ (thus define a fibered category over $\Hc(\Nop, T)$). Moreover, the functor
\[
C_\infty(\underline{t'}) \longrightarrow C_\infty(\underline{t})
\]
associated to a morphism of inverse systems
\[
\underline{t} \longrightarrow \underline{t'}
\]
is an \textit{isomorphism} in Hot. So the
\msp{page 118}
functor
\[
\underline{t} \longmapsto C_\infty(\underline{t}), \qquad \Hc(\Nop, T) \longrightarrow (\Hot)
\]
factors through the fundamental groupoid
\[
\Pi(\Hc(\Nop, T)) \longrightarrow (\Hot).
\]
When this groupoid, i.e. when the category $\Hc$ itself is connected, we can therefore conclude that all the $C_\infty(\underline{t})$ are isomorphic to one another in Hot. To conclude even that there is a \textit{transitive} system of isomorphisms between the $C_\infty(\underline{t})$, i.e. that $C_\infty(\underline{t})$ is ``unique up to \textit{unique} isomorphism'', one would have to know in addition that the fundamental groupoid is equivalent to the one-point category (trivial fundamental group), or, as one says, that $\Hc$ is \textit{simply connected,} or 1-connected.

The connectedness of $\Hc$ is assured when the category $T$ has (as is the case for $\underline{\mathrm{Ty}}$) a law of composition $(t,t') \longmapsto t * t'$, with morphisms functorial in $t, t'$
\[
\begin{tikzcd}[row sep=small]
 & t \\
t * t' \arrow[ru] \arrow[rd] & \\
 & t'
\end{tikzcd}
\]
(this last condition being automatic
% foot: Version 22 décembre 2001
\printedpage{75}
% head: A. Grothendieck : Dérivateurs, Ch. VII, Édition des Universités de Montpellier II et Paris VII
\msp{page 119}
if the law of composition admits the final object of $T$, assumed to exist, as a two-sided unit). In this case, $\underline{t} = (t_n)_n$ and $\underline{t'} = (t'_n)_n$ are linked to one another by $\underline{t} * \underline{t'} = (t_n * t'_n)_n$, since we have a morphism of inverse systems
\[
\begin{tikzcd}[column sep=small]
 & \underline{t} * \underline{t'} \arrow[ld] \arrow[rd] & \\
\underline{t} & & \underline{t'}.
\end{tikzcd}
\]
I do not know, for that matter, whether the arrow in Hot
\[
C_\infty(\underline{t}) \longrightarrow C_\infty(\underline{t'})
\]
associated to $\underline{t'} \longrightarrow \underline{t}$ in $\Hc$ depends only on $\underline{t}$, $\underline{t'}$, and not on the choice of $\underline{t'} \longrightarrow \underline{t}$. What is clear is that its image in $\Hotz$ is unique, it is the identity when we identify $C_\infty(\underline{t})$ and $C_\infty(\underline{t'})$ with one and the same element of $\Hotz$ defined by $\varinjlim H_T$.

And yet, I realize just now that, under the hypothesis of the existence of the law of composition and of its two-sided unit, and assuming in addition that the law in question is associative up to isomorphism, \textit{we have transitivity} in the isomorphisms (canonical!) between the $C_\infty(\underline{t})$ which we have just indicated!

\msp{page 120}
(Without having for all that to compute the fundamental group of $\Hc(\Nop, T)$!) If we have three elements $\underline{t}$, $\underline{t'}$, $\underline{t''}$ of $\Hc$, it suffices to look at the diagram deduced from the following commutative diagram
\[
\begin{tikzcd}[column sep=small]
 & & \underline{t} * \underline{t'} * \underline{t''} \arrow[ld] \arrow[rd] & & \\
 & \underline{t} * \underline{t'} \arrow[ld] \arrow[rd] & & \underline{t'} * \underline{t''} \arrow[ld] \arrow[rd] & \\
\underline{t} & & \underline{t'} & & \underline{t''}
\end{tikzcd}
\]
by applying to it the contrafunctor $\underline{s} \longrightarrow C_\infty(\underline{s})$ from $\Hc$ to Hot.

So in the case of the $C(a,b) = \underline{\mathrm{Ch}}_\Sigma(X; a, b)$, I no longer have any scruples about regarding the category $C_\infty(a,b)$ as a \textit{canonical} object of Hot, linked to the canonical ind-object $H_T(a,b)$ in Hot by a canonical homomorphism
\[
C_\infty(a,b) \longleftarrow H_T(a,b), \tag*{$(*)$}
\]
which in $\Hotz$ gives an isomorphism
\[
C_\infty(a,b) \simeq \varinjlim H_T(a,b), \tag*{$(**)$}
\]
i.e. the arrow $(*)$ induces, for $\xi \in \mathrm{Ob}\,\Hot$, a set-theoretic map
\[
\mathrm{Hom}_{\Hot}(\xi, C_\infty(a,b)) \longleftarrow \underbrace{\mathrm{Hom}_{\mathrm{Ind}\,\Hot}(\xi, H_T(a,b))}_{\underset{n}{\varinjlim}\, \mathrm{Hom}_{\Hot}(\xi, C_n(a,b))},
\]
which is \textit{bijective} if $\xi$ is in Hot.
% foot: Version 22 décembre 2001

\printedpage{76}
% head: A. Grothendieck : Dérivateurs, Ch. VII, Édition des Universités de Montpellier II et Paris VII
\msp{page 121}
Of course, the arrows
\[
C_\infty(\underline{t'}) \longrightarrow C_\infty(\underline{t}) \tag*{$(*)$}
\]
in Hot, associated to the arrows $\underline{t} \longrightarrow \underline{t'}$ in $\Hc(\Nop, T)$, are in fact arrows in % [sic] associés -> associées, accord avec « les flèches »
(Cat) itself - and even fully faithful functors, in the ``geometric'' case and for the two canonical arrows
\[
\begin{tikzcd}[column sep=small]
 & \underline{t} * \underline{t'} \arrow[ld] \arrow[rd] & \\
\underline{t} & & \underline{t'}.
\end{tikzcd}
\]
They are isomorphic \textit{in} Hot, because they define a homomorphism on the associated ind-objects of Cat\,hot, hence also for the associated ind-objects in Hot and in $\Hotz$, from which it follows that
\[
\mathrm{Hom}_{\Hot}(Y, C_\infty(\underline{t'})) \longrightarrow \mathrm{Hom}_{\Hot}(Y, C_\infty(\underline{t}))
\]
associated to $(*)$ is a bijection if $Y$ \textit{is finite} (\footnote[26]{Which suffices to ensure that $(*)$ is \texttt{[an]} isomorphism in Hot: $\Hot \longrightarrow \Hotz$ is conservative!}), this arrow then identifies with
\[
\underset{n}{\varinjlim}\, \mathrm{Hom}_{\Hot}(Y, C_{t'_n}) \longrightarrow \underset{n}{\varinjlim}\, \mathrm{Hom}_{\Hot}(Y, C_{t_n}).
\]
(For $Z \longmapsto \mathrm{Hom}_{\Hot}(Y, Z)$ commutes with filtering direct limits, if $Y$ is a category which is \textit{finite.} This comes from the fact that a filtering $\varinjlim$ of Thomason-Kan categories is likewise \dots) But I do not know, on the other hand, whether $(*)$ is \texttt{[an]} isomorphism in Cat\,hot. Only that for a \textit{finite} category $Y$, it induces
\[
\mathrm{Hom}_{\mathrm{Cat\,hot}}(Y, C_\infty(\underline{t'})) \xrightarrow{\ \sim\ } \mathrm{Hom}_{\mathrm{Cat\,hot}}(Y, C_\infty(\underline{t})),
\]
the two sides then identifying with
\msp{page 122}
\[
H_T(Y)_{\mathrm{Cat\,hot}} = \mathrm{Hom}_{\mathrm{Ind\,Cat\,hot}}(Y, H_T).
\]
To sum up these insights and perplexities about the $C_\infty(\underline{t})$:

\begin{enumerate}
\item[a)] To every $\underline{t} = (t_n)_{n \in \Nop} \in \Hc(\Nop, T)$ is associated a direct system, indexed by $\mathbf{N}$, of categories $C_n(\underline{t})$ $(= C_{t_n})$, and their direct limit
\[
C_\infty(\underline{t}).
\]
% foot: Version 22 décembre 2001

\printedpage{77}
% head: A. Grothendieck : Dérivateurs, Ch. VII, Édition des Universités de Montpellier II et Paris VII
\item[b)] This direct system, and also its limit is \texttt{[a]} contravariant functor in $\underline{t}$. If $\underline{t}, \underline{t'} \in \mathrm{Ob}\,\Hc$, then we can compare $C_\infty(\underline{t})$, $C_\infty(\underline{t'})$ and the corresponding direct systems, thanks to
\[
\begin{tikzcd}[column sep=small]
 & \underline{t} * \underline{t'} \arrow[ld] \arrow[rd] & \\
\underline{t} & & \underline{t'}.
\end{tikzcd}
\]

\item[c)] The object
\[
(C_n(\underline{t}))_{n \in \mathbf{N}} \in \mathrm{Ob\,Ind\,Cat\,hot}
\]
defined by the direct system $(C_n(\underline{t}))_{n \in \mathbf{N}}$ in Cat is (up to unique isomorphism) independent of the choice of $\underline{t}$, say $H_C$ or $H$. The homomorphisms
\[
C_\infty(\underline{t'}) \xrightarrow{u^*} C_\infty(\underline{t}) \qquad \text{associated to } u : \underline{t} \longrightarrow \underline{t'} \tag*{$(*)$}
\]
are compatible with this datum, i.e. we have commutativity in
\[
\begin{tikzcd}
\underset{\mathrm{Ind\,Cat\,hot}}{\varinjlim} C_*(\underline{t'}) \arrow[rr] & & \underset{\mathrm{Ind\,Cat\,hot}}{\varinjlim} C_*(\underline{t}) \\
 & H \arrow[lu, "\sim"] \arrow[ru, "\sim"'] &
\end{tikzcd}
\]

\msp{page 123}
\item[d)] It follows that the homomorphisms $(*)$ are \textit{isomorphisms} in Hot, hence all the $C_\infty(\underline{t})$ are isomorphic \textit{in} Hot. The trick with the $t * t'$ (in the case where we have a law $t * t'$ which is associative and admits the final object as unit object) even gives \textit{canonical} isomorphisms
\[
C_\infty(\underline{t'}) \longrightarrow C_\infty(\underline{t}) \qquad \text{in Hot,}
\]
and this is a transitive system of isomorphisms. So we have in Hot a canonical object
\[
C_\infty \in \mathrm{Ob}\,\Hot,
\]
whose relation to $H \in \mathrm{Ob\,Ind\,Hot}$ is deduced from c), it is a canonical arrow
\[
H \longrightarrow C_\infty \qquad \text{in Ind\,Hot (\footnotemark[27]).}
\]
\footnotetext[27]{\textbf{NB} $H$ is practically never essentially constant, so the arrow opposite \texttt{[is]} not an isomorphism in Ind\,Hot.}

This arrow induces an isomorphism for the corresponding objects of $\Hotz$; but viewing $H$ not as an ind-object of $\Hotz$, but taking the functor $H_0$ induced on $\Hot_0$ by the ind-representable functor $H$ on Hot. This induced functor is ``representable'' by the object $C_\infty$ of the larger category
\[
H_0 \simeq C_\infty \qquad \text{in } \Hotz.
\]

(If the canonical functor $\Hot \longrightarrow \Hotz$ is fully faithful, then this formula describes $C_\infty$ as an object of Hot, up to unique isomorphism, in terms of $H_0$, hence also in terms of $H \in \mathrm{Ob\,Ind\,Cat\,hot}$.)
% foot: Version 22 décembre 2001

\printedpage{78}
% head: A. Grothendieck : Dérivateurs, Ch. VII, Édition des Universités de Montpellier II et Paris VII
\msp{page 124}
\item[e)] I do not know, on the other hand, whether the arrows
\[
C_\infty(\underline{t'}) \longrightarrow C_\infty(\underline{t})
\]
are isomorphisms in Cat\,hot, although they are defined in terms of homomorphisms of direct systems of categories, which are isomorphisms in Ind\,Cat\,hot. This is so, even for the arrows $C_\infty(\underline{t}), C_\infty(\underline{t'}) \rightrightarrows C_\infty(\underline{t} * \underline{t'})$ used to compare $C_\infty(\underline{t})$ and $C_\infty(\underline{t'})$ canonically. Thus, I do not see how to define arrows in the reverse direction in Cat, which would be candidates to be inverse isomorphisms in Cat\,hot.

\item[f)] In the geometric case ($T = \underline{\mathrm{Ty}}$, situation described in terms of $X$, $\Sigma$, and $a, b \in \mathrm{Ob}\,X$, hence $C_\infty(\underline{t}; a, b)$, $C_\infty(\overline{a, b})$, $H(a,b)$ etc.), we \textit{do not have} a path composition operation at the level of the \textit{categories} $C_\infty(\underline{t}; a, b)$ and $C_\infty(\underline{t}; b, c)$, but only at the level of the objects $H(a,b)$, $H(b,c)$ of Ind\,Cat\,hot or of Ind\,Hot. Am I even sure I know how to define \textit{in} Hot
\[
C_\infty(a,b) \times C_\infty(b,c) \longrightarrow C_\infty(a,c) \ \ ?
\]

Surely, there are no problems for the images in $\Hotz$, since there they are the same as those of $H(a,b)$, $H(b,c)$, $H(a,c)$ respectively, and we have perfect, associative pairings, with unit, and with inverses where there ought to % [sic] là ou -> là où, accent omis sur le relatif
\msp{page 125}
be some. But not for the $C_\infty(a,b)$ themselves?
\end{enumerate}

It is a bit vexing that with all this work it would seem that we haven't even been able to damn well define the analogue in Cat of the classical loop space $\Omega(X, x)$ as a group \textit{in} Hot. (What does \textsc{Thomason} do?) But here is a thinkable idea [?]. To define \textit{all the same}
\[
C_\infty(a,b) \times C_\infty(b,c) \longrightarrow C_\infty(a,c)
\]
in Hot, we are going to describe the two functors $C_\infty(a,b)$ and $C_\infty(b,c)$ by different inverse systems $\underline{t}$, $\underline{t'}$. We can make it explicit by starting from a $\tau_0$ of useful length $\geq 2$, and taking in both cases
\[
t_n = t'_n = \tau_0^n,
\]
but different transition morphisms
\[
\begin{aligned}
(t_{n+p} \longrightarrow t_n) &= \pi_p * \mathrm{id}_{\tau_0^n} \\
(t'_{n+p} \longrightarrow t'_n) &= \mathrm{id}_{\tau_0^n} * \pi_p,
\end{aligned}
\]
where $\pi_p : \tau_0^p \longrightarrow 1$ is the total degeneracy. We will define $C_\infty(a,c)$ in terms of $\underline{t''} = \underline{t} * \underline{t'}$, hence
\[
\left\{
\begin{aligned}
t''_n &= \tau_0^{2n} \\
(t''_{n+p} \longrightarrow t''_n) &= \pi_p * \mathrm{id}_{\tau_0^{2n}} * \pi_p.
\end{aligned}
\right.
\]
% foot: Version 22 décembre 2001

\end{document}
